\documentclass{extreport}

\usepackage[14pt]{extsizes}
\usepackage[T2A]{fontenc}
\usepackage[utf8]{inputenc}
\usepackage[english,ukrainian]{babel}

\usepackage[a4paper, top=10mm, bottom=15mm, left=20mm, right=15mm]{geometry}
\usepackage{amsmath,amsfonts,amssymb,amsthm,mathtools}
\usepackage{listings}
\usepackage{graphicx}
\usepackage{enumitem}
\usepackage{verbatim}
\usepackage{listings}
\usepackage{xcolor}
\usepackage{textgreek}
\usepackage{diagbox}
\usepackage{pgfplots}

\pgfplotsset{compat = 1.16}


\lstdefinestyle{output_txt}{
    basicstyle=\ttfamily\footnotesize,
    breakatwhitespace=false,         
    breaklines=true,                 
    captionpos=b,                    
    keepspaces=true,                                    
    numbersep=5pt,                  
    showspaces=false,                
    showstringspaces=false,
    showtabs=false,                  
    tabsize=2
}
\definecolor{codegreen}{rgb}{0,0.6,0}
\definecolor{codegray}{rgb}{0.5,0.5,0.5}
\definecolor{codepurple}{rgb}{0.58,0,0.82}
\lstdefinestyle{python_code}{ 
    commentstyle=\color{codegreen},
    keywordstyle=\color{magenta},
    numberstyle=\tiny\color{codegray},
    stringstyle=\color{codepurple},
    basicstyle=\ttfamily\footnotesize,
    breakatwhitespace=false,         
    breaklines=true,                 
    captionpos=b,                    
    keepspaces=true,                            
    numbersep=5pt,                  
    showspaces=false,                
    showstringspaces=false,
    showtabs=false,                  
    tabsize=4
}

\setlist[enumerate]{nosep}
\graphicspath{{pics/}}
\DeclareGraphicsExtensions{.png}

\begin{document}
\begin{titlepage}
    \thispagestyle{empty}
    \begin{center}
        \includegraphics[width = \textwidth]{kpi}
        Міністерство освіти і науки України\\
        Національний технічний університет України\\
        <<Київський політехнічний інститут ім. І. Сікорського>>\\
        Інститут прикладного системного аналізу
    \end{center}
    \vspace{40mm}
    \begin{center}
        \textbf{Лабораторна робота} \\
        з курсу <<Методи оптимізації>> \\
        з теми <<Метод лінеаризації в нелінійному програмуванні>>
    \end{center}
    \vspace{20mm}
    \begin{flushleft}
        Виконали студенти 3 курсу групи КА-81 \\
        Галганов Олексій \\
        Єрко Андрій \\
        Фордуй Нікіта
    \end{flushleft}
    \begin{flushright}
        Перевірили \\
        Спекторський Ігор Якович \\
        Яковлева Алла Петрівна
    \end{flushright}
    \vspace{30mm}
    \begin{center}
        \textbf{Київ 2021}
    \end{center}
\end{titlepage}


\noindent\textbf{Завдання.} Реалізувати метод лінеаризації для задачі мінімізації нелінійної функції на множині, що задана набором нелінійних нерівностей:
\begin{gather}
    \begin{cases}
        f_0(x) \to \min \\
        f_i(x) \leq 0, \; i = 1,...,m
    \end{cases}
\end{gather}
Усі функції $f_0, f_1, ..., f_m \in C^1 (\mathbb{R}^n)$ --- неперервно диференційовні на $\mathbb{R}^n$.

\noindent\textbf{Опис методу.}
Покладемо $\Phi_N (x) = f(x) + N\cdot F(x)$ для великого $N>0$. $F(x) = \max\left\{0, f_1(x), ..., f_m(x)\right\} \geq 0$
для всіх $x \in \mathbb{R}^n$, причому рівність досягається тоді і тільки тоді, коли $f_i(x) \leq 0, \; i = 1,...,m$.
$N$ є одним з параметрів методу. Нехай $x_0 \in \mathbb{R}^n$ --- довільна початкова точка. На $k$-тому кроці роботи методу спочатку розв'язується
задача квадратичного програмування відносно $p$:
\begin{gather}\label{2}
    \begin{cases}
        \frac{1}{2} \Vert p \Vert^2 + \left<f'_0(x^k), p\right> \to \underset{p}{\min} \\
        f_i(x^k) + \left<f'_i(x^k), p\right> \leq 0, \; i = 1, ..., m
    \end{cases}
\end{gather}
Позначимо
\begin{gather*}
    D = \begin{pmatrix}
        -f'_1 (x^k) \\
        -f'_2 (x^k) \\
        \dots \\
        -f'_m(x^k)
    \end{pmatrix} = 
    \begin{pmatrix}
        -(\mathrm{grad}f_1 (x^k))^T \\
        -(\mathrm{grad}f_2 (x^k))^T \\
        \dots \\
        -(\mathrm{grad}f_m (x^k))^T
    \end{pmatrix},
    f = \begin{pmatrix}
        f_1(x^k) \\
        f_2(x^k) \\
        \dots \\
        f_m(x^k)
    \end{pmatrix},
    c = f'_0(x^k)
\end{gather*}
Тоді задача (\ref{2}) запишеться як
\begin{gather}\label{3}
    \begin{cases}
        \frac{1}{2} \Vert p \Vert^2 + \left<c, p\right> \to \underset{p}{\min} \\
        D p \geq f
    \end{cases}
\end{gather}
Ця задача є лінійним наближенням вихідної задачі в околі точки $x^k$.
Якщо $p^k$ --- розв'язок цієї задачі, то $x_{k+1} = x^k + \alpha^k p^k$, де $\alpha^k$ --- перше число в послідовності $1, \frac{1}{2}, \frac{1}{4}, ...$,
при якому виконується нерівність
\begin{gather}
    \Phi_N (x^k + \alpha^k p^k) \leq \Phi_N(x^k) - \varepsilon \alpha^k \Vert p^k \Vert^2
\end{gather}

Б. М. Пшеничний в книзі <<Метод лінеаризації>> \cite{1} доводить, що точка $x\in\mathbb{R}^n$ є стаціонарною для вихідної задачі тоді і тільки тоді, коли
задача квадратичного програмування
\begin{gather}
    \begin{cases}
        \frac{1}{2} \Vert p \Vert^2 + \left<f'_0(x), p\right> \to \underset{p}{\min} \\
        f_i(x) + \left<f'_i(x), p\right> \leq 0, \; i = 1, ..., m
    \end{cases}
\end{gather}
має розв'язок $p(x)$ і він рівний нулю. Там само доводиться, що за деяких умов (наприклад, ліпшицевість градієнтів функцій, що задають обмеження), метод є збіжним
до стаціонарних точок.

Квадратичну задачу (\ref{3}), що виникає на кожному кроці роботи методу, можна замінити так званою дуальною (двоїстою) задачею:
\begin{gather}
    \begin{cases}
        -\frac{1}{2} \Vert u \Vert^2 + \left<f, v\right> \to \underset{u, v}{\max} \\
        D^T v - u = c\\
        v \geq 0
    \end{cases} \Leftrightarrow
    \begin{cases}
        \frac{1}{2} \Vert u\Vert^2 - \left<f, v\right> \to \underset{u, v}{\min} \\
        D^T v - u = c \\
        v \geq 0
    \end{cases}
\end{gather}
Виключивши з задачі $u = D^T v - c$, отримаємо
\begin{gather}\label{non_neg}
    \begin{cases}
        \frac{1}{2}\left<D D^T v, v\right> - \left<D c + f, v \right> \to \underset{v}{\min} \\
        v \geq 0
    \end{cases}
\end{gather}
Якщо $v^*$ --- її розв'язок, то $p = u^* = D^T v^* - c$ буде розв'язком вихідної квадратичної задачі, як показано в \cite{1} та \cite{2}.

Задачу (\ref{non_neg}) можна розв'язати \emph{методом мультиплікативних оновлень}, викладеним у \cite{3}:
для задачі
\begin{gather*}
    \begin{cases}
        \frac{1}{2}\left<A v, v\right> + \left<b, v \right> \to \underset{v}{\min} \\
        v \geq 0
    \end{cases}
\end{gather*}
з $A > 0$ покладають $A^{+} = \max\left\{0, A\right\}$, $A^{-} = -\min\left\{0, A\right\}$ (поелементно).
Нехай $v^{0}$ --- деяке початкове наближення, на кожній ітерації його координати оновлюють за формулою
\begin{gather*}
    v^{k+1}_i = v^{k}_i \cdot \left(\frac{
        - b_i + \sqrt{b_i^2 + 4 \left(A^+ v^{k}\right)_i \left(A^- v^{k}\right)_i}
    }{
        2 \left(A^+ v^{k}\right)_i
    }\right), \; i = 1, ..., n
\end{gather*}
Ці оновлення можна записати як $v^{k+1} = \varphi\left(v^{k}\right)$. В \cite{3} доводиться, що таке
відображення $\phi$ має нерухому точку, яка є розв'язком задачі мінімізації, тому
послідовність $v^{k}$ збігається до розв'язку задачі.

Нарешті, у \cite{1} показано, що замість вибору параметру $N$ у методі лінеаризації
можна задавати деяке велике початкове значення $N_0$, а потім на кожній ітерації оновлювати його
за допомогою вектора $v^k$, отриманого при розв'язанні дуальної задачі (\ref{non_neg}), за правилом
\begin{gather*}
    \begin{cases}
        N_{k+1} = 2\sum\limits_{i=1}^m v^k_i, & \text{ якщо } \sum\limits_{i=1}^m v^k_i > N_k \\
        N_{k+1} = N_k, & \text{ інакше}
    \end{cases}
\end{gather*}

Також зауважимо, що метод можна застосувати до обмежень-рівностей виду $f(x) = 0$,
оскільки вони еквівалентні виконанню двох обмежень-нерівностей $f(x) \leq 0$ та $f(x) \geq 0$.

\newpage
\noindent\textbf{Результати роботи.} Ми перевірили роботу методу на декількох задачах зі ста випадкових початкових точок:
\begin{center}
    \begin{tabular}{|l|c|}
        \hline
        \textbf{задача} & \textbf{середня к-сть ітерацій} \\
        \hline
        $\begin{cases}
            x^2 + y^2 + z^2 \to \min \\
            x + y + z = 1
        \end{cases}$ & $4.17$ \\
        \hline
        $\begin{cases}
            x + 4y + z \to \min \\
            x^2 + 3y^2 + 2z^2 \leq 1
        \end{cases}$ & $241.5$ \\
        \hline
        $\begin{cases}
            (x-2)^2 + (y-1)^2 \to \min \\
            y - x^2 \leq 0
        \end{cases}$ & $20.23$ \\
        \hline 
        $\begin{cases}
            -xy \to \min \\
            x^2 + y^2 \leq 1 \\
            x \geq 0, y \geq 0
        \end{cases}$ & $20.9$ \\
        \hline
        $\begin{cases}
            x^2 + y \to \min \\
            x + y - 1 \leq 0 \\
            x^2 + y^2 \leq 9
        \end{cases}$ & $27.7$ \\
        \hline  
    \end{tabular}
\end{center}

В цілому, метод лінеаризації досить швидко збігається до розв'язку задачі,
але іноді збіжність була довгою через вибір початкової точки.

\noindent\textbf{Лістинг.}
Текст програми було розділено на \texttt{Optimizer.py} з реалізацією
методу лінеаризації та \texttt{lab5.py}, де викликаються необхідні функції та зберігаються результати.

\noindent\texttt{Optimizer.py}
\lstinputlisting[language = Python, style = python_code]{../code/Optimizer.py}

\noindent\texttt{lab4.py}
\lstinputlisting[language = Python, style = python_code]{../code/lab5.py}

\noindent\textbf{Висновки.} Реалізований нами метод лінеаризації виявився досить
швидким та універсальним методом розв'язання задач умовної мінімізації нелінійної
функції на множині, заданій нелінійними обмеженнями-нерівностями та рівностями.
Він є складним у реалізації та з обчислювальної точки зору, оскільки
на кожному кроці необхідно розв'язувати задачу квадратичного програмування з лінійними обмеженнями,
але ця складність виправдовується універсальністю застосування методу: цільова функція та функції-обмеження
мають бути лише неперервно-диференційовними. Але через те, що опуклість функцій-обмежень та строга опуклість
цільової функції не вимагається, в загальному випадку метод збігається лише до стаціонарних точок,
тому в таких випадках краще запустити мінімізацію з декількох початкових точок та порівняти отримані значення.

\begin{thebibliography}{3}
    \bibitem{1} Пшеничный Б.Н. \emph{Метод линеаризации}. --- М.: Наука. Главная редакция физико-математической литературы, 1983. --- 136 с.
    \bibitem{2} Dorn, W., 1960. Duality in quadratic programming. \emph{Quarterly of Applied Mathematics}, 18(2), pp.155-162.
    \bibitem{3} Sha, F., Lin, Y., Saul, L. and Lee, D., 2007. Multiplicative Updates for Nonnegative Quadratic Programming. Neural Computation, 19(8), pp.2004-2031.
\end{thebibliography}
\end{document}